\begin{exercisechunk}
\item \textbf{Width, dead ReLUs, and tangent coverage.}\suppmark
\label[exercise]{ex:l11-wide-relu}
\exerciseprofile{2}{2}{2}{\exproof\quad\excalculation\quad\exinterpretation}{%
Show how independent ReLU gates make complete gradient failure unlikely and
relate the resulting calculation to fixed-kernel convergence.}
For $w=(w_1,\ldots,w_m)\in\R^m$, let
\[
 f_m(w)=\frac1{\sqrt m}\sum_{r=1}^m[w_r]_+,
 \qquad
 L_m(w)=\frac12\bigl(f_m(w)-1\bigr)^2.
\]
Let $S_1,\ldots,S_m$ be independent random signs with
$\P(S_r=1)=\P(S_r=-1)=1/2$, fix $a\in(0,1)$, and initialize
\[
 W_{r,0}=\frac{aS_r}{\sqrt m}.
\]
Gradient descent uses $W_{t+1}=W_t-\eta\nabla L_m(W_t)$ with
$0<\eta\le1$. Define
\[
 Q=\sum_{r=1}^m\I\{S_r=1\}.
\]

\begin{enumerate}[(a),beginpenalty=10000]
\item Show that
\[
 \P(Q=0)=2^{-m},
\]
and that on the event $\{Q=0\}$ the iterates never move and
$L_m(W_t)=1/2$ for every $t\ge0$.

\item Condition on $Q=q>0$. Prove that the initially active coordinates
remain positive, the initially inactive coordinates remain negative, and
\begin{align*}
 1-f_m(W_t)
 &=\left(1-\eta\frac qm\right)^t
       \bigl(1-f_m(W_0)\bigr),\\
 L_m(W_t)
 &=\left(1-\eta\frac qm\right)^{2t}L_m(W_0).
\end{align*}

\item You may use Hoeffding's inequality in the form
\[
 \P\left(\frac Qm<\frac12-\delta\right)
 \le e^{-2m\delta^2},
 \qquad 0<\delta<\frac12.
\]
Prove that, with probability at least $1-e^{-2m\delta^2}$,
\[
 L_m(W_t)
 \le\left(1-\eta(1/2-\delta)\right)^{2t}L_m(W_0)
 \qquad\text{for every }t\ge0.
\]
Explain which failure of the single-ReLU model is removed by width and which
fact about the individual inactive units remains unchanged.

\item \emph{Optional extension.}
Let $K_\infty\in\R^{n\times n}$ be symmetric and positive semidefinite, and
suppose a residual vector follows
\[
 \dot r(t)=-K_\infty r(t).
\]
If
$K_\infty=U\operatorname{diag}(\lambda_1,\ldots,\lambda_n)U^\top$
with $U$ orthogonal, prove that
\[
 r(t)=U\operatorname{diag}
       (e^{-\lambda_1t},\ldots,e^{-\lambda_nt})U^\top r(0).
\]
Deduce that $\|r(t)\|_2^2/2\le
e^{-2\lambda t}\|r(0)\|_2^2/2$ when
$K_\infty\succeq\lambda I_n$. Explain both what the fixed-kernel assumption
makes exactly analyzable and what kind of network behavior it excludes.
\end{enumerate}

\begin{hints}
For (b), first compute $f_m(W_0)$ and the derivative with respect to one
active coordinate. Use induction to justify that the activation pattern does
not change. For the optional extension, express the differential equation in
the eigenbasis of $K_\infty$.
\end{hints}
\end{exercisechunk}
